% Source: https://artofproblemsolving.com/community/c3h2949432_2022_emcc_guts_round_p13_p24__exeter_math_club_competition
% Stable identity: aops-topic-2949432-q15
A regular tetrahedron of edge $\sqrt6$ has height $\sqrt6\cdot\sqrt{\dfrac23}=2$. So with $z$ the height above $ABC$, $P$ is at $z=2$ on the axis and $Q$ at $z=1$. At height $z$, the cross-section of $ABCP$ is a copy of $ABC$ scaled by $1-\dfrac z2$, and the cross-section of $A'B'C'Q$ is a copy with the same orientation scaled by $\dfrac{z-1}2$, both centred on the axis.
\begin{asy}
usepackage("amsmath");
size(5cm);
defaultpen(fontsize(10pt));
draw((0,0)--(6,0)--(6,3)--(0,3)--cycle);pair P=(3,2),Q=(3,1);
draw((0,0)--P--(6,0),dashed);draw((0,3)--Q--(6,3),dashed);
fill((2.25,1.5)--P--(3.75,1.5)--Q--cycle,lightgray);draw((2.25,1.5)--P--(3.75,1.5)--Q--cycle);
dot(P);dot(Q);label("$P$",P,N);label("$Q$",Q,S);label("$3$",(6,1.5),E);label("$z=\dfrac32$",(3.75,1.5),E);
\end{asy}

The intersection at height $z$ is the smaller triangle. The two scales are equal at $z=\dfrac32$, both $\dfrac14$. So the intersection is a double pyramid with apexes $Q$ and $P$ over a triangle of area $\dfrac1{16}\cdot\dfrac{\sqrt3}4\cdot6=\dfrac{3\sqrt3}{32}$, total height $1$:
\[V=\dfrac13\cdot\dfrac{3\sqrt3}{32}\cdot1=\dfrac{\sqrt3}{32}.\]
So $m+n=3+32=\boxed{35}$.
